\chapter{Jordan--Wigner Transformation}

\section{Definition}

The Jordan--Wigner transformation maps a one-dimensional spin-$1/2$ chain to a
chain of spinless fermions. The essential difficulty is that spin operators on
different sites commute, whereas fermionic operators on different sites must
anticommute. The nonlocal Jordan--Wigner string supplies precisely the required
fermionic signs.

Throughout this chapter we use the convention
\begin{equation}
    |0\rangle_j \equiv |\uparrow\rangle_j,
    \qquad
    |1\rangle_j \equiv |\downarrow\rangle_j .
\end{equation}
Thus the occupied fermionic state corresponds to spin down. The Pauli matrices are
\begin{equation}
    \sigma^x =
    \begin{pmatrix}
        0 & 1\\
        1 & 0
    \end{pmatrix},
    \qquad
    \sigma^y =
    \begin{pmatrix}
        0 & -\ii\\
        \ii & 0
    \end{pmatrix},
    \qquad
    \sigma^z =
    \begin{pmatrix}
        1 & 0\\
        0 & -1
    \end{pmatrix}.
\end{equation}
The spin raising and lowering operators are
\begin{equation}
    \sigma^+
    =
    \frac{\sigma^x+\ii\sigma^y}{2}
    =
    \begin{pmatrix}
        0 & 1\\
        0 & 0
    \end{pmatrix},
    \qquad
    \sigma^-
    =
    \frac{\sigma^x-\ii\sigma^y}{2}
    =
    \begin{pmatrix}
        0 & 0\\
        1 & 0
    \end{pmatrix}.
\end{equation}
They act as
\begin{equation}
    \sigma^+|\downarrow\rangle = |\uparrow\rangle,
    \qquad
    \sigma^-|\uparrow\rangle = |\downarrow\rangle .
\end{equation}

With the convention that spin down is occupied, it is natural to identify the local
hard-core boson operators as
\begin{equation}
    b_j=\sigma_j^+,
    \qquad
    b_j^\dagger=\sigma_j^- .
\end{equation}
The local number operator is then
\begin{equation}
    n_j
    =
    b_j^\dagger b_j
    =
    \sigma_j^-\sigma_j^+
    =
    \frac{1-\sigma_j^z}{2}.
\end{equation}
Equivalently,
\begin{equation}
    \sigma_j^z = 1-2n_j.
\end{equation}

We define the Jordan--Wigner string by
\begin{equation}
\label{eq:jw-string-definition}
    K_j
    =
    \prod_{\ell=1}^{j-1}\sigma_\ell^z
    =
    \prod_{\ell=1}^{j-1}(1-2n_\ell).
\end{equation}
Since \(n_\ell\) has eigenvalues \(0\) and \(1\), we may also write
\begin{equation}
    K_j
    =
    \exp\!\left(
        \ii\pi\sum_{\ell=1}^{j-1}n_\ell
    \right).
\end{equation}
This exponential form is often useful because
\begin{equation}
    \ee^{\ii\pi n_\ell}=1-2n_\ell=\sigma_\ell^z.
\end{equation}

The Jordan--Wigner transformation used in these notes is
\begin{equation}
\label{eq:jw-definition-product}
    c_j
    =
    \left(\prod_{\ell=1}^{j-1}\sigma_\ell^z\right)\sigma_j^+,
    \qquad
    c_j^\dagger
    =
    \left(\prod_{\ell=1}^{j-1}\sigma_\ell^z\right)\sigma_j^- .
\end{equation}
Equivalently,
\begin{equation}
\label{eq:jw-definition-compact}
    c_j = K_j\sigma_j^+ = K_j b_j,
    \qquad
    c_j^\dagger = K_j\sigma_j^- = K_j b_j^\dagger .
\end{equation}
The inverse transformation is
\begin{equation}
\label{eq:jw-inverse}
    \sigma_j^+ = K_j c_j,
    \qquad
    \sigma_j^- = K_j c_j^\dagger,
\end{equation}
or, equivalently,
\begin{equation}
    b_j=K_jc_j,
    \qquad
    b_j^\dagger=K_jc_j^\dagger .
\end{equation}

The string satisfies
\begin{equation}
    K_j^\dagger=K_j,
    \qquad
    K_j^2=1,
    \qquad
    K_{j+1}=K_j\sigma_j^z=K_j(1-2n_j).
\end{equation}
Because \(K_j\) only acts on sites strictly to the left of \(j\), it commutes with all
operators on site \(j\):
\begin{equation}
    [K_j,\sigma_j^\pm]=0,
    \qquad
    [K_j,\sigma_j^z]=0,
    \qquad
    [K_j,c_j]=0,
    \qquad
    [K_j,c_j^\dagger]=0.
\end{equation}

Using the inverse transformation, the spin operators become
\begin{equation}
\label{eq:spin-operators-jw}
    \sigma_j^x
    =
    K_j(c_j^\dagger+c_j),
    \qquad
    \sigma_j^y
    =
    K_j\,\ii(c_j^\dagger-c_j),
    \qquad
    \sigma_j^z
    =
    1-2c_j^\dagger c_j .
\end{equation}

\begin{remark}
A common source of sign confusion is the alternative string
\(\prod_{\ell<j}(-\sigma_\ell^z)\). This differs from the convention used here by
a site-dependent sign:
\begin{equation}
    \prod_{\ell=1}^{j-1}(-\sigma_\ell^z)
    =
    (-1)^{j-1}
    \prod_{\ell=1}^{j-1}\sigma_\ell^z .
\end{equation}
It therefore corresponds to a staggered fermion redefinition
\begin{equation}
    c_j^{\mathrm{alt}} = (-1)^{j-1}c_j .
\end{equation}
In this chapter and in the rest of these notes, we consistently use
\(K_j=\prod_{\ell<j}\sigma_\ell^z\).
\end{remark}

\section{Canonical anticommutation relations}

We now prove that the operators defined by the Jordan--Wigner transformation obey
the canonical fermionic anticommutation relations
\begin{equation}
\label{eq:canonical-fermion-algebra}
    \{c_i,c_j^\dagger\}=\delta_{ij},
    \qquad
    \{c_i,c_j\}=0,
    \qquad
    \{c_i^\dagger,c_j^\dagger\}=0.
\end{equation}

First consider the same site. Since \(K_j\) commutes with \(\sigma_j^\pm\) and
\(K_j^2=1\),
\begin{align}
    c_jc_j^\dagger
    &=
    K_j\sigma_j^+K_j\sigma_j^-
    =
    \sigma_j^+\sigma_j^-,
    \\
    c_j^\dagger c_j
    &=
    K_j\sigma_j^-K_j\sigma_j^+
    =
    \sigma_j^-\sigma_j^+.
\end{align}
Therefore
\begin{equation}
    \{c_j,c_j^\dagger\}
    =
    \sigma_j^+\sigma_j^-
    +
    \sigma_j^-\sigma_j^+
    =
    1.
\end{equation}
Also,
\begin{equation}
    c_j^2=(\sigma_j^+)^2=0,
    \qquad
    (c_j^\dagger)^2=(\sigma_j^-)^2=0.
\end{equation}

Now take \(i<j\). Define
\begin{equation}
    R_{ij}
    =
    \prod_{\ell=i+1}^{j-1}\sigma_\ell^z,
\end{equation}
so that
\begin{equation}
    K_j = K_i\sigma_i^zR_{ij}.
\end{equation}
All operators at different spin sites commute. The only nontrivial sign comes from
the factor \(\sigma_i^z\), which is present in the string of the operator at site \(j\).

For two annihilation operators,
\begin{align}
    c_i c_j
    &=
    K_i\sigma_i^+K_j\sigma_j^+
    \nonumber\\
    &=
    K_i\sigma_i^+K_i\sigma_i^zR_{ij}\sigma_j^+
    \nonumber\\
    &=
    \sigma_i^+\sigma_i^zR_{ij}\sigma_j^+
    \nonumber\\
    &=
    -\sigma_i^+R_{ij}\sigma_j^+,
\end{align}
where we used
\begin{equation}
    \sigma_i^+\sigma_i^z=-\sigma_i^+.
\end{equation}
On the other hand,
\begin{align}
    c_j c_i
    &=
    K_j\sigma_j^+K_i\sigma_i^+
    \nonumber\\
    &=
    K_i\sigma_i^zR_{ij}\sigma_j^+K_i\sigma_i^+
    \nonumber\\
    &=
    \sigma_i^z\sigma_i^+R_{ij}\sigma_j^+
    \nonumber\\
    &=
    +\sigma_i^+R_{ij}\sigma_j^+,
\end{align}
where
\begin{equation}
    \sigma_i^z\sigma_i^+=+\sigma_i^+.
\end{equation}
Hence
\begin{equation}
    c_ic_j=-c_jc_i,
    \qquad i\neq j.
\end{equation}
Thus
\begin{equation}
    \{c_i,c_j\}=0.
\end{equation}

Similarly,
\begin{align}
    c_i c_j^\dagger
    &=
    K_i\sigma_i^+K_j\sigma_j^-
    \nonumber\\
    &=
    \sigma_i^+\sigma_i^zR_{ij}\sigma_j^-
    \nonumber\\
    &=
    -\sigma_i^+R_{ij}\sigma_j^-,
\end{align}
while
\begin{align}
    c_j^\dagger c_i
    &=
    K_j\sigma_j^-K_i\sigma_i^+
    \nonumber\\
    &=
    \sigma_i^z\sigma_i^+R_{ij}\sigma_j^-
    \nonumber\\
    &=
    +\sigma_i^+R_{ij}\sigma_j^-.
\end{align}
Therefore
\begin{equation}
    c_ic_j^\dagger=-c_j^\dagger c_i,
    \qquad i\neq j.
\end{equation}
Together with the same-site result, this proves
\begin{equation}
    \{c_i,c_j^\dagger\}=\delta_{ij}.
\end{equation}

Finally, for two creation operators,
\begin{align}
    c_i^\dagger c_j^\dagger
    &=
    K_i\sigma_i^-K_j\sigma_j^-
    \nonumber\\
    &=
    \sigma_i^-\sigma_i^zR_{ij}\sigma_j^-
    \nonumber\\
    &=
    +\sigma_i^-R_{ij}\sigma_j^-,
\end{align}
where
\begin{equation}
    \sigma_i^-\sigma_i^z=+\sigma_i^-,
\end{equation}
whereas
\begin{align}
    c_j^\dagger c_i^\dagger
    &=
    K_j\sigma_j^-K_i\sigma_i^-
    \nonumber\\
    &=
    \sigma_i^z\sigma_i^-R_{ij}\sigma_j^-
    \nonumber\\
    &=
    -\sigma_i^-R_{ij}\sigma_j^-.
\end{align}
Hence
\begin{equation}
    c_i^\dagger c_j^\dagger=-c_j^\dagger c_i^\dagger,
    \qquad i\neq j.
\end{equation}
This completes the proof of the canonical fermionic anticommutation relations.

\section{Boundary conditions and fermion parity}

For an open chain, the Jordan--Wigner transformation is unambiguous: the string
always extends from the left boundary to the site immediately before \(j\). No operator
wraps around the chain.

For a periodic spin chain, however, a boundary term connecting site \(L\) to site \(1\)
carries a fermion-parity-dependent sign. The total fermion number and fermion parity
are
\begin{equation}
    N=\sum_{j=1}^L n_j,
    \qquad
    P_F=(-1)^N.
\end{equation}
Since \(1-2n_j=\sigma_j^z\), we have
\begin{equation}
\label{eq:fermion-parity-spin}
    P_F
    =
    \prod_{j=1}^L(1-2n_j)
    =
    \prod_{j=1}^L\sigma_j^z.
\end{equation}
The parity operator satisfies
\begin{equation}
    P_F^2=1,
    \qquad
    P_F c_j P_F=-c_j,
    \qquad
    P_F c_j^\dagger P_F=-c_j^\dagger.
\end{equation}

The formal Jordan--Wigner string at site \(L+1\) is
\begin{equation}
    K_{L+1}
    =
    \prod_{j=1}^L\sigma_j^z
    =
    P_F.
\end{equation}
Thus, if the spin variables are periodic,
\begin{equation}
    \sigma_{L+1}^\pm\equiv\sigma_1^\pm,
\end{equation}
the formal extension of the Jordan--Wigner definition gives
\begin{equation}
\label{eq:formal-boundary-fermion}
    c_{L+1}^{\mathrm{form}}
    =
    K_{L+1}\sigma_{L+1}^+
    =
    P_F\sigma_1^+
    =
    P_Fc_1 .
\end{equation}
This equation is a formal extension of the operator definition. It should not yet
be identified with the effective boundary condition of a quadratic fermion
Hamiltonian. The effective boundary condition is obtained only after the boundary
spin bilinears are rewritten in the same algebraic form as the bulk nearest-neighbour
bilinears.

For example,
\begin{align}
    b_L^\dagger b_1
    &=
    K_Lc_L^\dagger c_1
    \nonumber\\
    &=
    \left(\prod_{\ell=1}^{L-1}\sigma_\ell^z\right)c_L^\dagger c_1
    \nonumber\\
    &=
    P_F\sigma_L^z c_L^\dagger c_1
    \nonumber\\
    &=
    -P_F c_L^\dagger c_1.
\end{align}
Similarly,
\begin{equation}
    b_L^\dagger b_1^\dagger
    =
    -P_F c_L^\dagger c_1^\dagger.
\end{equation}
Therefore, within a fixed fermion-parity sector \(P_F=p=\pm1\), rewriting the boundary
term in the same form as the bulk terms is equivalent to imposing the effective
fermionic boundary condition
\begin{equation}
\label{eq:effective-jw-boundary-condition}
    c_{L+1}=-p\,c_1.
\end{equation}
Consequently,
\begin{equation}
    p=+1
    \quad\Longrightarrow\quad
    c_{L+1}=-c_1,
\end{equation}
so the even-parity sector corresponds to anti-periodic fermions, whereas
\begin{equation}
    p=-1
    \quad\Longrightarrow\quad
    c_{L+1}=c_1,
\end{equation}
so the odd-parity sector corresponds to periodic fermions.

\section{Useful Jordan--Wigner identities}

This section collects the identities most frequently used when translating
one-dimensional spin-chain Hamiltonians into fermionic language. All formulae below
use the convention
\begin{equation}
    K_j=\prod_{\ell=1}^{j-1}\sigma_\ell^z
    =
    \prod_{\ell=1}^{j-1}(1-2n_\ell),
    \qquad
    c_j=K_j\sigma_j^+,
    \qquad
    c_j^\dagger=K_j\sigma_j^-.
\end{equation}

\subsection*{Local spin, hard-core boson, and number operators}

The hard-core boson convention is
\begin{equation}
    \sigma_j^+=b_j,
    \qquad
    \sigma_j^-=b_j^\dagger,
    \qquad
    n_j=b_j^\dagger b_j=c_j^\dagger c_j.
\end{equation}
Therefore
\begin{equation}
    \sigma_j^x=b_j^\dagger+b_j,
    \qquad
    \sigma_j^y=\ii(b_j^\dagger-b_j),
    \qquad
    \sigma_j^z=1-2b_j^\dagger b_j=1-2c_j^\dagger c_j.
\end{equation}
The local hard-core algebra is
\begin{equation}
    \{b_j,b_j^\dagger\}=1,
    \qquad
    \{b_j,b_j\}=0,
    \qquad
    \{b_j^\dagger,b_j^\dagger\}=0,
\end{equation}
while for \(i\neq j\),
\begin{equation}
    [b_i,b_j]=0,
    \qquad
    [b_i,b_j^\dagger]=0,
    \qquad
    [b_i^\dagger,b_j^\dagger]=0.
\end{equation}
The fermions satisfy
\begin{equation}
    \{c_i,c_j^\dagger\}=\delta_{ij},
    \qquad
    \{c_i,c_j\}=0,
    \qquad
    \{c_i^\dagger,c_j^\dagger\}=0.
\end{equation}

\subsection*{String identities}

The string identities most often used are
\begin{equation}
    K_j^2=1,
    \qquad
    K_j^\dagger=K_j,
    \qquad
    K_{j+1}=K_j(1-2n_j).
\end{equation}
Moreover,
\begin{equation}
    c_j^\dagger(1-2n_j)=c_j^\dagger,
    \qquad
    c_j(1-2n_j)=-c_j,
\end{equation}
and
\begin{equation}
    (1-2n_j)c_j^\dagger=-c_j^\dagger,
    \qquad
    (1-2n_j)c_j=c_j.
\end{equation}
These four identities are the main source of the signs in nearest-neighbour JW
formulae.

\subsection*{Hard-core boson bilinears}

Using \(b_j=K_jc_j\), \(b_j^\dagger=K_jc_j^\dagger\), and
\(K_{j+1}=K_j(1-2n_j)\), one obtains
\begin{align}
    b_j^\dagger b_j
    &=
    c_j^\dagger c_j
    =
    n_j,
    \\
    b_j^\dagger b_{j+1}^\dagger
    &=
    c_j^\dagger(1-2n_j)c_{j+1}^\dagger
    =
    c_j^\dagger c_{j+1}^\dagger,
    \\
    b_j^\dagger b_{j+1}
    &=
    c_j^\dagger(1-2n_j)c_{j+1}
    =
    c_j^\dagger c_{j+1},
    \\
    b_j b_{j+1}
    &=
    c_j(1-2n_j)c_{j+1}
    =
    -c_jc_{j+1},
    \\
    b_j b_{j+1}^\dagger
    &=
    c_j(1-2n_j)c_{j+1}^\dagger
    =
    -c_jc_{j+1}^\dagger.
\end{align}
Equivalently, after reordering fermions on different sites,
\begin{equation}
    -c_jc_{j+1}^\dagger
    =
    c_{j+1}^\dagger c_j,
    \qquad
    -c_jc_{j+1}
    =
    c_{j+1}c_j.
\end{equation}

\subsection*{Spin operators in fermions}

The single-site spin operators are
\begin{equation}
    \sigma_j^x
    =
    K_j(c_j^\dagger+c_j),
    \qquad
    \sigma_j^y
    =
    K_j\,\ii(c_j^\dagger-c_j),
    \qquad
    \sigma_j^z
    =
    1-2n_j.
\end{equation}

The nearest-neighbour \(xx\) bilinear is
\begin{align}
    \sigma_j^x\sigma_{j+1}^x
    &=
    (b_j^\dagger+b_j)(b_{j+1}^\dagger+b_{j+1})
    \nonumber\\
    &=
    c_j^\dagger c_{j+1}^\dagger
    +
    c_j^\dagger c_{j+1}
    +
    c_{j+1}^\dagger c_j
    +
    c_{j+1}c_j
    \nonumber\\
    &=
    c_j^\dagger c_{j+1}^\dagger
    +
    c_j^\dagger c_{j+1}
    +
    \text{H.c.}.
\end{align}
The nearest-neighbour \(yy\) bilinear is
\begin{align}
    \sigma_j^y\sigma_{j+1}^y
    &=
    - (b_j^\dagger-b_j)(b_{j+1}^\dagger-b_{j+1})
    \nonumber\\
    &=
    -c_j^\dagger c_{j+1}^\dagger
    +
    c_j^\dagger c_{j+1}
    +
    c_{j+1}^\dagger c_j
    -
    c_{j+1}c_j
    \nonumber\\
    &=
    -\left(
        c_j^\dagger c_{j+1}^\dagger
        -
        c_j^\dagger c_{j+1}
        +
        \text{H.c.}
    \right).
\end{align}
The nearest-neighbour \(zz\) bilinear is
\begin{align}
    \sigma_j^z\sigma_{j+1}^z
    &=
    (1-2n_j)(1-2n_{j+1})
    \nonumber\\
    &=
    1-2n_j-2n_{j+1}+4n_jn_{j+1}.
\end{align}

Useful linear combinations are
\begin{align}
    \sigma_j^x\sigma_{j+1}^x
    +
    \sigma_j^y\sigma_{j+1}^y
    &=
    2\left(
        c_j^\dagger c_{j+1}
        +
        c_{j+1}^\dagger c_j
    \right),
    \\
    \sigma_j^x\sigma_{j+1}^x
    -
    \sigma_j^y\sigma_{j+1}^y
    &=
    2\left(
        c_j^\dagger c_{j+1}^\dagger
        +
        c_{j+1}c_j
    \right).
\end{align}

\subsection*{Tensor-product representation}

For direct matrix construction on a finite chain of length \(L\), the JW operators are
\begin{equation}
    c_j
    =
    \left(
        \bigotimes_{\ell=1}^{j-1}\sigma_\ell^z
    \right)
    \otimes
    \sigma_j^+
    \otimes
    \left(
        \bigotimes_{\ell=j+1}^{L}\id_\ell
    \right),
\end{equation}
and
\begin{equation}
    c_j^\dagger
    =
    \left(
        \bigotimes_{\ell=1}^{j-1}\sigma_\ell^z
    \right)
    \otimes
    \sigma_j^-
    \otimes
    \left(
        \bigotimes_{\ell=j+1}^{L}\id_\ell
    \right).
\end{equation}
For example, for \(L=4\),
\begin{align}
    c_1
    &=
    \sigma_1^+\otimes\id_2\otimes\id_3\otimes\id_4,
    \\
    c_2
    &=
    \sigma_1^z\otimes\sigma_2^+\otimes\id_3\otimes\id_4,
    \\
    c_3
    &=
    \sigma_1^z\otimes\sigma_2^z\otimes\sigma_3^+\otimes\id_4,
    \\
    c_4
    &=
    \sigma_1^z\otimes\sigma_2^z\otimes\sigma_3^z\otimes\sigma_4^+.
\end{align}
The creation operators are obtained by replacing \(\sigma_j^+\) with \(\sigma_j^-\).

\subsection*{Compact summary}

The essential dictionary is
\begin{equation}
    n_j
    =
    c_j^\dagger c_j
    =
    \frac{1-\sigma_j^z}{2},
    \qquad
    \sigma_j^z=1-2n_j,
\end{equation}
\begin{equation}
    K_j
    =
    \prod_{\ell=1}^{j-1}\sigma_\ell^z
    =
    \prod_{\ell=1}^{j-1}(1-2n_\ell),
\end{equation}
\begin{equation}
    c_j=K_j\sigma_j^+,
    \qquad
    c_j^\dagger=K_j\sigma_j^-,
\end{equation}
\begin{equation}
    \sigma_j^x=K_j(c_j^\dagger+c_j),
    \qquad
    \sigma_j^y=K_j\,\ii(c_j^\dagger-c_j),
    \qquad
    \sigma_j^z=1-2c_j^\dagger c_j,
\end{equation}
and
\begin{equation}
    \{c_i,c_j^\dagger\}=\delta_{ij},
    \qquad
    \{c_i,c_j\}=0,
    \qquad
    \{c_i^\dagger,c_j^\dagger\}=0.
\end{equation}
These identities are the basic toolkit for translating one-dimensional spin chains
into fermionic language.